\section{Experimental Results and Comparative Analysis}
\label{sec:experimental_results}

To evaluate the efficiency and robustness of our proposed approaches, we conducted an empirical comparison between the exact integer programming formulation (\textit{FEASIBILITY} solver) and the proposed \textit{SMART HEURISTIC}. The experimental dataset originally encompasses 449 problem instances derived from both SOP (Sequential Ordering Problem) and generic generated network structures. However, a significant portion of these instances proved intractable for the exact solver due to time and memory limitations. While the \textit{SMART HEURISTIC} successfully solved all 449 instances, providing robust upper bounds across the entire dataset, the exact \textit{FEASIBILITY} solver was only able to converge to an optimal solution for 187 instances. The comparative analysis below primarily focuses on the subset of 187 commonly solved instances to directly compare execution time (runtime in seconds) and solution quality (measured by the number of jumps), followed by a discussion on the instances solved exclusively by the heuristic.

\subsection{Summary of Results by Instance Group}

The table below presents a summarized view of the experimental results grouped by the structural size and type of the instances. It details the number of instances in each group, their average topological characteristics (nodes, edges, and density), and a comparison of the average number of jumps and runtimes for both the exact solver (\textit{FEASIBILITY}) and the \textit{SMART HEURISTIC}. Note that for groups where the exact solver experienced intractability, its averages only reflect the subset of instances it successfully solved, whereas the heuristic's averages reflect all instances. Groups marked with a timeout denote complete intractability for the exact solver under the given time and memory constraints.

\begin{table}[hbt!]
\centering
\caption{Average characteristics, solution quality, and runtime grouped by instance category.}
\label{tab:summary_results}
\resizebox{\textwidth}{!}{
\begin{tabular}{l c c c c | c c | c c}
\toprule
\multirow{2}{*}{\textbf{Group}} & \multirow{2}{*}{\textbf{\# Inst.}} & \multicolumn{3}{c}{\textbf{Average Characteristics}} & \multicolumn{2}{|c}{\textbf{Average Jumps}} & \multicolumn{2}{|c}{\textbf{Average Runtime (s)}} \\
\cmidrule(lr){3-5} \cmidrule(lr){6-7} \cmidrule(lr){8-9}
& & \textbf{$|V|$} & \textbf{$|E|$} & \textbf{Density} & \textbf{Exact} & \textbf{Heuristic} & \textbf{Exact} & \textbf{Heuristic} \\
\midrule
n10 & 58 & 13.0 & 18.2 & 0.81 & 3.52 & 3.60 & 0.56 & 0.0031 \\
n20 & 60 & 23.0 & 42.1 & 0.76 & 7.12 & 8.02 & 240.06 & 0.0107 \\
n30 & 60 & 33.0 & 74.5 & 0.73 & 10.45 & 13.02 & 3116.82 & 0.0313 \\
n40 & 50 & 43.0 & 124.4 & 0.67 & - & 19.66 & Timeout & 0.0805 \\
n50 & 60 & 53.0 & 164.7 & 0.71 & 17.53 & 23.72 & 11769.89 & 0.1380 \\
n60 & 60 & 63.0 & 219.6 & 0.71 & - & 29.22 & Timeout & 0.2380 \\
n70 & 60 & 73.0 & 280.1 & 0.70 & - & 34.73 & Timeout & 0.3995 \\
sop & 41 & 94.5 & 443.7 & 0.35 & 2.06 & 43.29 & 20.54 & 40.3613 \\
\bottomrule
\end{tabular}
}
\end{table}

\subsection{Computational Performance and Solution Quality}

The experimentation reveals a contrast in the computational effort required by the two algorithms. The exact solver exhibits an average execution time of $1,931.47$ seconds across all instances, with high variance; median solving time is roughly $9.82$ seconds, but worst-case instances consumed up to $27,861$ seconds (approximately $7.7$ hours). In contrast, the heuristic framework recorded instantaneous execution times (measured at $0.0$ seconds within the timer's precision threshold) across the entire dataset.

Despite this astronomical reduction in computational effort, the smart heuristic maintains highly competitive solution quality. It successfully matched the exact optimal number of jumps in $67.38\%$ of the evaluated instances. In cases where the heuristic failed to find the absolute optimal solution ($32.62\%$ of cases), the optimality gap was notably small, deviating by an average of only $1.44$ jumps. Over the entire dataset, the mean jump deviation between the heuristic and the exact solver is merely $0.47$ jumps. For the specific subset of SOP instances, the heuristic achieved a $100\%$ exact match rate, suggesting extreme reliability on structurally defined problems.

\subsection{Impact of Instance Characteristics}

To better understand when the problem becomes intractable for the exact solver and when the heuristic begins to deviate, we performed a Pearson correlation analysis against key topological metrics of the underlying instances: the number of nodes ($|V|$), the number of edges ($|E|$), the mean in-degree, and the closed edge density.

The execution time of the exact solver is predominantly governed by the network's density and size. Specifically, runtime displays a strong positive correlation with the number of edges ($r = 0.71$) and the mean in-degree ($r = 0.50$). Interestingly, the raw number of nodes exhibits a moderate correlation ($r = 0.44$), indicating that the complexity of solving the Arboreal Jump Number problem exactly is dictated more by the interconnectedness of the graph (edges) than by its raw vertex count.

Similarly, the heuristic's solution quality is influenced by these same density metrics. The optimality gap (difference in jumps) shows a significant positive correlation with the mean in-degree ($r = 0.66$) and the number of edges ($r = 0.57$). Highly interconnected components appear to obfuscate the heuristic's local decision-making process, slightly increasing the likelihood of suboptimal jump assignments.

\subsection{Intractability and Heuristic Exclusivity}

The most compelling argument for the adoption of the \textit{SMART HEURISTIC} is its ability to find valid, high-quality jump configurations for instances where the exact formulation completely fails to converge within reasonable limits. Out of the 449 evaluated problem instances, 262 instances (approximately $58\%$) were solved exclusively by the heuristic. 

These 262 intractable instances are heavily skewed towards larger and more complex networks, further highlighting the exponential scaling bottlenecks of the exact integer programming approach. Specifically, the heuristic-only solutions comprise 60 instances from the $n=70$ generated group, 60 instances from the $n=60$ group, 50 instances from the $n=40$ group, and 41 instances from the $n=50$ group, alongside 23 large SOP instances. As the graph size ($|V| \ge 40$) and interconnectivity increase, the exact solver's search tree becomes prohibitively large. In these scenarios, the \textit{SMART HEURISTIC} transitions from being merely a fast alternative to being the only viable computational tool for the Arboreal Jump Number problem, delivering instantaneous results where exact mathematical formulation hits a computational wall.

\subsection{Algorithm Selection Strategy}

Based on the empirical evidence, the selection between the exact solver and the smart heuristic should be dictated by the structural parameters of the input graph, primarily utilizing the number of edges ($|E|$) as a decision threshold:

\begin{itemize}
    \item \textbf{Small to Medium, Sparse Instances ($|E| \le 100$):} The exact solver is highly recommended. Within this subset, the mean runtime drops to approximately $1,117$ seconds (roughly $18$ minutes), a generally acceptable timeframe for offline computation. In this region, the heuristic achieves a $70.06\%$ exact match rate. If absolute mathematical guarantees are required, the exact solver's computational cost is justifiable.
    
    \item \textbf{Large, Dense Instances ($|E| > 100$):} The smart heuristic is strictly recommended. For instances exceeding $100$ edges, the exact solver suffers from exponential time degradation, averaging $16,343$ seconds (over $4.5$ hours). Meanwhile, the smart heuristic continues to deliver instantaneous results with a highly acceptable error margin (an average deviation of $1.40$ jumps). 
\end{itemize}

In conclusion, the \textit{SMART HEURISTIC} presents a highly effective operational compromise, breaking the exponential bottleneck of the exact integer formulation while sacrificing only minimal solution quality on complex, dense instances.